% paper_fresh_beats.tex -- a from-scratch "simple, unedited" top-beats version of the paper (2026-10-28), written from
% the author's decisions D1-D20 (writeup/INTEGRATION_PLAN_2026-10-26.md) and the generated number files only; the
% current paper text was not read. Build from writeup/fresh: pdflatex, bibtex, pdflatex, pdflatex (the \input@path
% and \graphicspath lines below resolve generated/ and figures/ one level up). Preamble copied from
% writeup/paper_rehaul.tex (lines before \begin{document}); additions are marked "fresh".
\documentclass[11pt]{article}
\usepackage[margin=1in]{geometry}
\usepackage[T1]{fontenc}
\usepackage{lmodern}
\usepackage{amsmath}
\usepackage{amssymb}
\usepackage{bm}
\DeclareMathOperator{\QLIKE}{QLIKE}
\DeclareMathOperator{\sign}{sign}
\DeclareMathOperator{\clip}{clip}
\DeclareMathOperator{\softplus}{softplus}
\DeclareMathOperator*{\argmin}{arg\,min}
% 2026-10-15 (notation pass): shared symbols, see writeup/NOTATION_GUIDE_2026-10-15.md
\newcommand{\RV}{\mathrm{RV}}
\newcommand{\E}{\mathbb{E}}
\DeclareMathOperator{\Var}{Var}
\DeclareMathOperator{\Cov}{Cov}
\newcommand{\SR}{\mathrm{SR}}
\newcommand{\DM}{\mathrm{DM}}
\usepackage{booktabs}
\usepackage{enumitem}
\usepackage[round]{natbib}
\usepackage{graphicx}
\usepackage{float}
\usepackage{caption}
\usepackage{threeparttable}
\usepackage{longtable}
% 2026-10-15 (table legibility): landscape pages for wide tables, stacked headers. siunitx is NOT loaded: the
% local MiKTeX's expl3 is too old for it (fatal "Support package expl3 too old"); use r columns with fixed decimals.
\usepackage{pdflscape}
\usepackage{makecell}
\usepackage{xcolor}
% 2026-10-28 (author, Table tab:rh_channels_abl legibility): \cellcolor shading in tables. colortbl v1.0k (2024)
% calls \insert@pcolumn, which array.sty defines only from v2.6; the local MiKTeX's array v2.5g builds p columns with
% \insert@column, so it stands in there (no effect where array is current, as on Overleaf).
\usepackage{colortbl}
\makeatletter\@ifundefined{insert@pcolumn}{\let\insert@pcolumn\insert@column}{}\makeatother
\captionsetup{font=small,labelfont=bf}
\usepackage{hyperref}
\emergencystretch=3em
\setlist{itemsep=1pt, topsep=3pt, parsep=0pt}
\setlist[itemize,2]{label=--}
\newcommand{\todo}[1]{\textit{[#1]}}
% 2026-10-26 (integration, JFEc): the author-notes switch. \authornotesfalse, the default, is the submission build:
% every \editc box prints nothing and the notes to the authors (rehaul/99_author_notes.tex) are not read.
% \authornotestrue restores the draft view (the boxes, the notes section and the draft subtitle).
\newif\ifauthornotes
\authornotesfalse
% 2026-10-16 (user): editorial notes for the consistency pass. Colored boxes in the draft, not part of the paper.
% red: two live statements disagree. orange: the statement does not match the computation behind a printed number.
% blue: a reader cannot recover the input or the rule. Numbered E1--E25 so a note can be named in editing.
% 2026-10-26 (integration, JFEc): the box is defined only under \authornotestrue; otherwise \editc takes its three
% arguments and prints nothing.
\ifauthornotes
\newcommand{\editc}[3]{%
  \par\smallskip\noindent
  \begingroup
  \setlength{\fboxsep}{3pt}%
  \colorbox{#1!12}{%
    \parbox{\dimexpr\linewidth-2\fboxsep\relax}{%
      \raggedright\footnotesize\textcolor{#1!75!black}{\textbf{#2.} #3}%
    }%
  }%
  \endgroup
  \par\smallskip
}
\else
\newcommand{\editc}[3]{\ignorespaces}
\fi
% 2026-10-26 (integration, JFEc): \rhPending marks a number whose generated macro is not written yet. \rhPending{}
% (the generators' form, and \rhPending alone) prints "pending"; \rhPending{name} prints the name of the missing macro,
% so the gap is visible in the PDF. Defined before the number files: their \providecommand{\rhPending}{...} lines
% leave this definition in place.
\makeatletter
\DeclareRobustCommand{\rhPending}{\@ifnextchar\bgroup{\rh@pending}{\textit{pending}}}
\newcommand{\rh@pending}[1]{%
  \if\relax\detokenize{#1}\relax\textit{pending}\else\textit{[pending: \texttt{\detokenize{#1}}]}\fi}
\makeatother
% 2026-10-19 (user): the draft names each exogenous series by its raw parquet column name for now (the "% data field"
% names of generated/rehaul_dm_dictionary.tex). \rawseriesfalse restores the symbols of the 15 October notation pass
% (writeup/NOTATION_GUIDE_ADDENDA_2026-10-15.md, A.2). \ser{raw}{symbol}: the series' name; \serx{raw}{raw
% superscript/subscript}{symbol}: for a symbol whose own superscript takes the horizon, e.g.
% \serx{sumret2_ewstock}{^{(25)}}{\mathrm{RV}^{\mathrm{EW},(25)}}. In text mode the raw name may break after an
% underscore (table cells); in math mode it is one box.
\newif\ifrawseries
\rawseriestrue
\makeatletter
\begingroup\catcode`\_=12
\gdef\ser@split#1_#2\@nil{#1\if\relax\detokenize{#2}\relax\else\textunderscore\allowbreak\ser@split#2\@nil\fi}
\endgroup
\def\ser@text#1{{\normalfont\ttfamily\expandafter\ser@split\detokenize{#1_}\@nil}}
\DeclareRobustCommand{\sername}[1]{\ifmmode\text{\ser@text{#1}}\else\ser@text{#1}\fi}
\makeatother
\newcommand{\ser}[2]{\ifrawseries\sername{#1}\else #2\fi}
\newcommand{\serx}[3]{\ifrawseries\sername{#1}#2\else #3\fi}
% fresh (2026-10-28): this file lives in writeup/fresh; generated/ and figures/ are one level up. \input@path makes
% every \input{generated/...} (here and inside the generated floats) resolve to ../generated/...; \graphicspath does
% the same for \includegraphics{figures/...} and \includegraphics{generated/...}.
\makeatletter\def\input@path{{../}}\makeatother
\graphicspath{{../}}
% fresh (2026-10-28): two product names that contain digits, as macros, so that the running text holds no digit.
\newcommand{\SPfh}{S\&P~500}
\newcommand{\zdte}{0DTE}
% 2026-10-15 (consolidation): every macro-only number file is loaded once, here. Each file holds definitions only (no
% tables, checked), and no macro name is defined in two of them. The in-body \input lines of these files in the section
% files are parked, so nothing is defined twice. Generators: make_table_close_main.py (close_main_numbers),
% make_fig_qlike_sharpe.py (fig_qlike_sharpe_numbers), make_rehaul_results.py, make_rehaul_economics.py,
% make_rehaul_fi.py, make_rehaul_appendix.py (rehaul_*_numbers).
% 2026-10-15 (number macros): added make_rehaul_data_methods.py (rehaul_dm_numbers: \rhdm... sample counts and
% constants), make_buckets_tex.py (buckets_numbers: \bkSize... group sizes), and the --numbers-only modes of
% experiments/close_{murphy_splits,compounding,shift_test}_headline.py (close_*_numbers: \hmurRhoESLossMid, \cmp...,
% \shiftSwing... / \shiftNoiseBand), whose table-mixed appendix_close_*.tex files are not loaded here.
\input{generated/close_main_numbers}
\input{generated/fig_qlike_sharpe_numbers}
\input{generated/rehaul_results_numbers}
\input{generated/rehaul_econ_numbers}
\input{generated/rehaul_fi_numbers}
\input{generated/rehaul_app_numbers}
\input{generated/rehaul_dm_numbers}
\input{generated/buckets_numbers}
\input{generated/close_murphy_splits_numbers}
\input{generated/close_compounding_numbers}
\input{generated/close_shift_test_numbers}
\input{generated/rehaul_oos_numbers}
% 2026-10-20 (user): Table 1 from the pooled linear reruns, like for like (writeup/make_rehaul_oos_rerun.py).
\input{generated/rehaul_oos_rerun_numbers}
% 2026-10-23 (user, final rerun): the pooled shift test and compounding (writeup/make_rehaul_pooled_extras.py: \rhps...,
% \rhpc...; notes E26, E27, N15b), the window robustness (writeup/make_rehaul_robust_window.py: \rhrb..., Table
% tab:rh_robust_w3500) and the pooled variable importance (writeup/make_rehaul_fi_pooled.py: \rhfip...; 08_feature_importance
% loads it only if it is not loaded here). Macro-only files; no name is defined in another loaded file (checked).
\input{generated/rehaul_pooled_extras_numbers}
\input{generated/rehaul_robust_w3500_numbers}
\input{generated/rehaul_fip_numbers}
% 2026-10-23 (coordinator, tree learning curve): the tree ensembles across training windows in Section 8
% (writeup/make_rehaul_fip_lc.py: \rhlc..., Table tab:rh_fip_lc). Macro-only file; no name is defined in another
% loaded file (checked).
\input{generated/rehaul_fip_lc_numbers}
% 2026-10-26 (user): the late-starting series' TreeSHAP shares in the years they are live, Section 8
% (writeup/make_rehaul_fi_live.py: \rhfil..., Figure fig:rh_fi_live). Macro-only file; no name is defined in another
% loaded file (checked).
\input{generated/rehaul_fi_live_numbers}
% 2026-10-25 (user): numbers reported by Maisch, Zhang, Ulrich and Zimmer (SSRN 7269778), for the comparison in the
% conclusion (writeup/make_external_numbers.py: \extMaisch...). Macro-only file; not results of ours; no name is
% defined in another loaded file (checked).
\input{generated/external_maisch_numbers}
% 2026-10-25 (user, step 1 of the gap analysis): the reach of the late-start rule (writeup/make_rehaul_live_windows.py:
% \rhlv...). Macro-only file; no name is defined in another loaded file (checked).
\input{generated/rehaul_live_numbers}
% 2026-10-26 (integration, JFEc): the macro files of the drafts integrated in this pass (D5, D6, D10; the ladder grid).
% Each is macro-only (every non-comment line is a definition, checked) and no macro name in them is defined in any
% other file loaded here (checked by script, 26 October 2026). Their floats load them only when undefined
% (\ifdefined ... \input), so do not \input them again in a section file: rehaul_ablation_numbers uses \newcommand,
% and a second load stops the build. Definition styles differ between the generated files: \newcommand (a duplicate
% stops the build), \providecommand alone (the first load wins) and \providecommand{}\renewcommand (the last load
% wins), so a new file is checked for name collisions, not for load order.
%   rehaul_ablation_numbers      writeup/make_rehaul_ablations.py        \rhab...  (tab:rh_ablation; drafts/2026-10-25_ablations)
%   rehaul_econ_main_numbers     writeup/make_rehaul_econ_main.py        \rhem...  (tab:rh_econ_main; drafts/2026-10-25_econ_main)
%   rehaul_econ_decomp_numbers   writeup/make_rehaul_econ_decomp.py      \rhed...  (drafts/2026-10-25_econ_decomp)
%   rehaul_mcs_numbers           writeup/make_rehaul_mcs.py              \rhmcs... (tab:rh_mcs_spa; drafts/2026-10-25_mcs)
%   rehaul_close_accuracy_numbers writeup/make_rehaul_close_accuracy.py  \rhca...  (drafts/2026-10-25_close_accuracy)
%   rehaul_econ_beta_numbers     writeup/make_rehaul_econ_beta.py        \rheb...  (tab:rh_econ_beta; drafts/2026-10-26_econ_beta)
%   rehaul_train_arms_numbers    writeup/make_rehaul_train_arms.py       \rhta...  (tab:rh_train_arms; drafts/2026-10-26_train_arms)
%   rehaul_complexity_numbers    writeup/make_rehaul_complexity.py       \rhcx...  (drafts/2026-10-26_complexity; pre-patch)
\input{generated/rehaul_ablation_numbers}
\input{generated/rehaul_econ_main_numbers}
\input{generated/rehaul_econ_decomp_numbers}
\input{generated/rehaul_mcs_numbers}
\input{generated/rehaul_close_accuracy_numbers}
\input{generated/rehaul_econ_beta_numbers}
\input{generated/rehaul_train_arms_numbers}
\input{generated/rehaul_complexity_numbers}
% ---- 2026-10-26 (integration, JFEc): RESERVED BLOCK for the NUMBERS agent's new macro files -------------------
% To come: the feature-count ladder on the FOMC-patched design (D13), the tuning disclosure (D14), the NYSE-holiday
% disclosure (D15) and Table 1 against the RTH13 OLS-HAR benchmark (D2). Add each here as a macro-only \input, after
% checking that none of its macro names is defined in a file loaded above. The front matter marks the numbers these
% files are to supply with \rhPending{name} (see the integration report and rehaul/99_author_notes.tex); replace
% each \rhPending{name} by the generated macro once the file is loaded. File names below are placeholders.
% 2026-10-26 (coordinator): the NUMBERS agent's files, now generated (no name collision with the files above, checked
% by the NUMBERS agent); the \rhPending{name} placeholders were replaced by these macros.
\input{generated/rehaul_ladder_numbers}
\input{generated/rehaul_tuning_numbers}
\input{generated/rehaul_holiday_numbers}
\input{generated/rehaul_oos_vs_ols_numbers}
\input{generated/rehaul_window_sweep_numbers}
\input{generated/rehaul_ci_har_numbers}
% 2026-10-28 (author: main2.tex on Overleaf was over its 2 MB limit for editable files): the number files that the
% sections load through \ifdefined guards are loaded here too, so flatten.py drops those guards instead of inlining
% the file at each one.
\input{generated/rehaul_channels_numbers}
\input{generated/rehaul_factor_numbers}
\input{generated/rehaul_beyond_ladder_numbers}
\input{generated/rehaul_beyond_factor_numbers}
% 2026-10-28 (author, alternative-data edit): \rhbyFLNSeries, the 36 series of the six groups beyond HAR + RM + FOMC.
\input{generated/rehaul_beyond_factor_lin_numbers}
% fresh (2026-10-28): the ends of the feeds and the data fixes (writeup/make_rehaul_disclosures.py: \rhds...), used by the
% Data section here. Macro-only file; no name is defined in another loaded file (checked against every generated file).
\input{generated/rehaul_disclosure_numbers}
% ---- end of the reserved block ------------------------------------------------------------------------------------

\title{Exogenous Predictors in Pooled Intraday Forecasts of Realized Variance\\\large Fresh beats version (simple,
unedited), written from the decisions and the results only}
\author{James Chen \and Mihai Cucuringu \and Christopher S. Jones \and Austin Pollok}
\date{Working draft, not for citation}

\begin{document}
\maketitle

\begin{abstract}
\noindent We forecast the realized variance of every regular-hours half hour of \SPfh{} futures with pooled penalized
linear models and tree ensembles. Against a pooled OLS-HAR benchmark, \rhdmNSeriesPooled{} exogenous series lower
QLIKE by \rhorVsOlsRedLinRng\% (linear) and \rhorVsOlsRedTreeRng\% (trees) but raise MSE, through a few crash days. The
futures' realized moments and the FOMC calendar carry \rhabBaseShareMin--\rhabBaseShareMax\% of the gain. Beyond them,
other series add no significant accuracy to linear models and a small, dense gain to boosted trees. Signing
\zdte{} SPX straddles with these forecasts raises the Sharpe ratio over the benchmark's, significantly for
\rhemFullVsOlsSigWord{} model of \rhmcsNAltWord.
\end{abstract}

\section*{Governing claim}

Pooled over the \rhorNBars{} regular-hours half hours of the day, forecasts of \SPfh{} futures realized variance that
use \rhdmNSeriesPooled{} exogenous series have a QLIKE \rhorVsOlsRedLinRng\% (penalized linear models) and
\rhorVsOlsRedTreeRng\% (tree ensembles) below the pooled OLS-HAR benchmark, at the cost of a higher MSE that a few days
of the spring crash produce. Between \rhabBaseShareMin\% and \rhabBaseShareMax\% of each class's gain comes from two
sources that are already studied: the futures' own realized moments on ordinary days and the FOMC calendar on statement
days. Beyond them, the remaining \rhbyFLNSeries{} series add no significant accuracy to penalized linear models and
change the boosted trees' QLIKE by \rhabLgbmFullVsAddRmFomcAllPct\% and \rhabXgbFullVsAddRmFomcAllPct\%, a gain that a
few principal or partial-least-squares components do not recover. The full-set forecasts sign a \zdte{} straddle with
a higher Sharpe ratio than the benchmark's in \rhemFullVsOlsPos{} of \rhmcsNAlt{} classes, by \rhemFullVsOlsMin{} to
\rhemFullVsOlsMax, significantly for \rhemFullVsOlsSigWord.

% ======================================================================================================================
\section{Introduction}
\label{sec:intro}

% CONFLICT: D18 makes the question what the other exogenous series add beyond HAR + RM + FOMC, while D1, D2 and D7 put
% the full-set-versus-OLS-HAR gain and the MSE clause in the abstract and introduction. These beats lead with the
% full-set gain (D1/D2/D7) and answer the D18 question second (contribution two); the author may prefer the reverse.
\begin{enumerate}
\item This paper asks what exogenous series add to a pooled forecast of intraday realized variance beyond the HAR
  components, the futures' own realized moments and the FOMC calendar.
\item The target is the realized variance of each regular-hours half hour of the E-mini \SPfh{} futures, the
  \rhorNBars{} bars ending \rhfipBarFirst{} to \rhfipBarLast{} ET, and one model of each class is fitted to the bars of
  every time of day together.
\item The benchmark is the OLS on the HAR benchmark set \citep{Corsi2009}, pooled on the same bars.
\item With the full set of \rhdmNSeriesPooled{} exogenous series (\rhorNColsAf{} predictors), the forecasts have a
  lower QLIKE than the benchmark in \rhorVsOlsFullLowerSig{} of \rhorVsOlsFullN{} model classes, by
  \rhorVsOlsRedLinRng\% for ridge, lasso and elastic net and \rhorVsOlsRedTreeRng\% for LightGBM, XGBoost and the random
  forest, on \rhorDays{} trading days from \rhorSpanStart{} to \rhorSpanEnd.
\item The same full-set forecasts have a higher MSE than their HAR-set fits; in \rhcaYearsOos{} the increase comes from
  \rhcaTopDays{} days between \rhcaTopFirstBarOos{} and \rhcaTopLastBarOos, and without the crash window
  \rhcaNMseLowerSigBarOosExCrash{} of \rhorVsOlsFullN{} classes have a significantly lower MSE on the full set.
\item The first contribution splits the gain into two channels: the \rhchNEventDays{} scheduled FOMC statement days
  hold \rhchEventBarPct\% of the scored bars but \rhchLinShareEventMin--\rhchLinShareEventMax\% of the linear models'
  gain and \rhchTreeShareEventMin--\rhchTreeShareEventMax\% of the trees' gain, and on those days the FOMC calendar
  alone delivers the gain.
\item On the other, ordinary days the futures' realized moments carry most of the gain, \rhchOrdRmShareTreeRng\% of it
  for the trees.
\item The second contribution answers the question: the HAR set plus the realized moments and the FOMC calendar
  (\rhbyFLNBase{} predictors) reaches \rhabBaseShareMin--\rhabBaseShareMax\% of each class's full-set gain, and the
  remaining \rhbyFLNSeries{} series add no significant accuracy to the penalized linear models.
\item The boosted trees gain a little more from those series, \rhabLgbmFullVsAddRmFomcAllPct\% (LightGBM) and
  \rhabXgbFullVsAddRmFomcAllPct\% (XGBoost) in QLIKE, and the gain is spread over the series: up to $k = 3$ principal or
  partial-least-squares components of them recover at most \rhbyFXgbBestShortOrdR{} of its ordinary-day part.
\item The third contribution is the application: a \zdte{} SPX straddle opened half an hour before the close and
  signed by the pooled forecast of the closing bar has a Sharpe ratio higher by \rhemFullVsOlsMin{} to
  \rhemFullVsOlsMax{} with full-set forecasts than with the OLS-HAR forecast.
\item That Sharpe ratio difference is significant for \rhemFullVsOlsSigWord{} of the \rhmcsNAltWord{} classes, and the
  profits come from a few long positions.
\item Related work, realized variance: the paper builds on realized variance as a measure of integrated variance
  \citep{AndersenBollerslevDieboldLabys2003,BarndorffNielsenShephard2002} and on the HAR model and its extensions with
  semivariances, quarticity and leverage \citep{Corsi2009,PattonSheppard2015,BollerslevPattonQuaedvlieg2016,CorsiReno2012}.
\item Related work, intraday forecasting: it pools the intraday bars of the regular session with predictors that see
  the overnight session, as in \citet{EngleSokalska2012} and \citet{ZhangZhangCucuringuQian2024}, and uses the
  overnight information studied by \citet{StroudJohannes2014}, \citet{ChenYuZivot2012} and \citet{TodorovaSoucek2014}.
\item Related work, machine learning: penalized regressions and tree ensembles have been applied to realized variance
  \citep{AudrinoKnaus2016,LuongDokuchaev2018,RahimikiaPoon2020,christensen2023machine}, and
  \citet{BrancoRubesamZevallos2024} ask whether anything beats linear models; here the answer depends on the predictor
  set.
\item Related work, alternative data: sentiment, attention and news measures have been tried as volatility predictors
  \citep{AudrinoSigristBallinari2020,BehrendtSchmidt2018,CaporinPoli2017,MittnikRobinzonovSpindler2015}; here they are
  read on the windows in which they are live.
\item Related work, economic evaluation: variance forecasts are valued in option strategies
  \citep{FlemingKirbyOstdiek2001,BandiFusariReno2023,AlmeidaFreireHizmeri2025,Vilkov2026}, and
  \citet{MaischZhangUlrichZimmer2026} trade machine-learning portfolios of half-hour SPX options in
  \extMaischYear.
\item Section~\ref{sec:rh_data} describes the data, Section~\ref{sec:rh_predictors} the predictors,
  Section~\ref{sec:rh_methods} the design and tests, Section~\ref{sec:rh_results} the accuracy results,
  Section~\ref{sec:rh_beyond} what the remaining series add, Section~\ref{sec:rh_econ} the straddle application and
  Section~\ref{sec:concl} concludes.
\end{enumerate}

% ======================================================================================================================
\section{Data}
\label{sec:rh_data}

\begin{enumerate}
\item The realized variance of a half hour is the sum of the squared one-minute log returns of the front-month E-mini
  \SPfh{} futures in it, on a grid of \rhdmNIntervals{} half-hour intervals a day that includes the overnight session.
\item The grid holds \rhdmGridRows{} intervals with an observed realized variance, from \rhPending{first day of the
  futures sample} to \rhdsSampleEnd.
\item The forecasts are scored on the \rhorNBars{} regular-hours bars of the \rhorDays{} trading days from
  \rhorSpanStart{} to \rhorSpanEnd, \rhorNIntB{} bars after dropping \rhorEarlyCloses{} early-close days and bars with a
  non-positive realized variance.
\item The exogenous series come from the index constituents (equal- and value-weighted realized moments, turnover,
  order flow and spreads), StockTwits (attention and sentiment), Cboe (the VIX indices and net volatility demand) and
  the FOMC calendar, each recorded at the end of its half-hour interval.
\item Several feeds end before the futures sample: the vendor's VIX prints stop on \rhdsVixVendorLast{} and the last
  \rhdsVixSpliceDays{} trading days are spliced from a public feed that agrees with the vendor on \rhdsVixAgree\% of
  the overlapping stamps, and the volatility-demand series stop on \rhdsVdLast.
\item The FOMC calendar is the patched one: it includes the scheduled meeting of \rhdsFomcSchedMeeting{} that the
  emergency action of \rhdsFomcReplaced{} replaced, and the proximity series are not winsorized; the span holds
  \rhchNEventDays{} scheduled statement days.
% NOTE (D15): before the official version, rerun every model with DROP_HOLIDAY_SESSIONS on.
\item The \rhholNSessions{} shortened futures sessions on NYSE holidays in the span stay in the estimation windows and
  are forecast (\rhholNBars{} bars) but not scored.
\item Dropping those sessions from training shrinks the linear models' every-bar QLIKE gain by \rhholEffDBBarPpRng{}
  percentage points, and the gain stays significant (Appendix~\ref{app:robust}).
\item The option data are SPXW quotes on expiration days: \rheNdays{} option days from \rheSampleStart{} to
  \rheSampleEnd, early closes dropped.
\end{enumerate}

% ======================================================================================================================
\section{Predictors}
\label{sec:rh_predictors}

\begin{enumerate}
\item The \rhdmNSeriesPooled{} exogenous series fall into \bkNBuckets{} groups: the futures' realized moments
  (\bkSizeMom{} series), liquidity (\bkSizeLiq), equal-weighted (\bkSizeMew) and value-weighted (\bkSizeMvw) market
  moments of the constituents, sentiment (\bkSizeSent), implied volatility (\bkSizeIv), volatility demand (\bkSizeVd)
  and the FOMC calendar (\bkSizeFomc).
\item The realized moments are the futures' return sum, absolute variation, realized skewness and quarticity, positive
  semivariance, bipower variation and first-order autocovariance; the FOMC group is the statement-day indicator and the
  inverse distances to the next and to the last statement.
\item Each series is divided by its trailing same-interval mean or standard deviation, the diurnal adjustment
  \citep{AndersenBollerslev1997,AndersenSuTodorovZhang2024}, then transformed and winsorized.
\item Each series enters as its average over the previous $k \in \{1, 5, 25, 125, 625, 3125\}$ half-hour intervals and
  as the share of those intervals on which it was observed, \rhfipNPerSeries{} predictors a series.
\item The averages run over the full grid of \rhdmNIntervals{} intervals a day, so the forecasts for the morning bars
  see the overnight session.
\item The HAR benchmark set holds the HAR components of the target at the same \rhdmNHorizonsPooled{} horizons and
  \rhorNCalendarAf{} calendar columns; the full set adds every exogenous series, \rhorNColsAf{} predictors in all.
\item The base set of Section~\ref{sec:rh_beyond}, the HAR set plus the realized moments and the FOMC calendar, has
  \rhbyFLNBase{} predictors; the remaining \rhbyFLNSeries{} series are the liquidity, market-moment, sentiment,
  implied-volatility and volatility-demand groups.
\item A series that starts after the futures data enters as zero until the rolling window holds \rhorWindowDays{} days
  of its history (the late-start rule).
\item Under that rule the VIX averages are live from \rhlvVixLiveYear, the StockTwits sentiment from
  \rhlvStSentLiveYear{} and the volatility demand from \rhlvVdLiveYear, on \rhlvVixEvRange\%, \rhlvStSentEvRange\% and
  \rhlvVdEvRange\% of the scored bars.
\end{enumerate}

% ======================================================================================================================
\section{Methods}
\label{sec:rh_methods}\label{sec:rh_methods_tests}\label{sec:rh_results_eval}

\begin{enumerate}
\item Every model is estimated on the \rhorNBars{} regular-hours bars only, pooled, on a rolling window of the previous
  \rhorWindowBars{} bars (\rhorWindowDays{} trading days), and re-estimated at every bar.
\item Each model forecasts a diurnally scaled volatility $\hat y_{t,i}$, and the variance forecast
  $\hat\sigma^2_{t,i} = \hat y^{2}_{t,i}\, s^2_{t,i}$ is scored against the realized variance of the bar.
\item The linear classes are OLS, ridge \citep{HoerlKennard1970}, lasso \citep{Tibshirani1996} and elastic net
  \citep{ZouHastie2005}; the tree classes are LightGBM \citep{Ke2017}, XGBoost \citep{Chen2016} and the random forest
  \citep{Breiman2001}.
\item The penalties are chosen once on \rhorPenSample{} and frozen; each tree method keeps one configuration from an
  Optuna search \citep{Akiba2019} on the \rhtuTreeSearchSpan{} pre-sample, also frozen.
\item The tuning only has to put off-the-shelf models in the forecasting regime, and its checks are disclosed in
  Appendix~\ref{sec:rh_methods_hyper}.
\item The primary benchmark is the OLS on the HAR set; the ridge, lasso and elastic net on the HAR set coincide with it
  (largest QLIKE change \rhorVsOlsPctLinHarMaxAbs\%).
\item The classical HAR extensions, SHAR, HARQ and LHAR, serve as augmented benchmarks
  \citep{PattonSheppard2015,BollerslevPattonQuaedvlieg2016,CorsiReno2012}.
\item QLIKE is the primary loss because it ranks variance forecasts consistently under a noisy proxy
  \citep{Patton2011}; MSE and the out-of-sample $R^2$ against the historical same-interval mean
  \citep{CampbellThompson2008} are reported with it.
\item Forecasts are compared with the Diebold--Mariano test \citep{DieboldMariano1995} with the
  \citet{HarveyLeybourneNewbold1997} correction, on the daily average loss over the day's bars, Newey--West lag
  \rhorHacLag.
\item The day types are fixed by the Fed's published calendar before any scoring: statement days are the scheduled FOMC
  statement days, ordinary days are the others, and analyses added after the first scoring are labelled as such.
\end{enumerate}

% ======================================================================================================================
\section{Forecast accuracy}
\label{sec:rh_results}

\input{generated/rehaul_results_oos}

% CONFLICT: D2 asks that every main-text forecast, the trees' full sets included, be tested against the RTH13 OLS-HAR.
% The generated Table 1 (tab:rh_oos) tests each tree's full set against its own HAR-set fit, and the generated table
% that tests every forecast against the OLS-HAR (tab:rh_oos_vsols) also carries 16:00-bar columns, which D16 keeps out of
% the main text. The beats below give the every-bar tests against the OLS-HAR from the \rhorVsOls... macros; an
% every-bar-only version of tab:rh_oos_vsols would resolve this.
\begin{enumerate}
\item The full-set forecasts beat the OLS-HAR benchmark in \rhorVsOlsFullLowerSig{} of \rhorVsOlsFullN{} classes
  (Table~\ref{tab:rh_oos}): QLIKE falls by \rhorVsOlsRedLinRng\% for the linear models and \rhorVsOlsRedTreeRng\% for the
  trees, with the largest Diebold--Mariano $p$-value \rhorVsOlsPMax.
\item Within each class the full set also beats the class's own HAR-set fit, with Diebold--Mariano statistics from
  \rhorDmRidgeB{} (ridge) to \rhorDmLgbmB{} (LightGBM).
\item The trees on the HAR set alone are within \rhorVsOlsPctTreeHarMaxAbs\% of the OLS-HAR in QLIKE, so the trees'
  gain comes from the exogenous series, not from a nonlinear use of the HAR components.
\item The trees beat the linear models: on QLIKE the model confidence set at both levels reported keeps only
  \rhmcsBarQlikeTmaxSetNinety{} of the \rhmcsNModelsWord{} distinct forecasts.
\item Among the linear full sets the lasso has the lowest QLIKE, \rhldPctLassoRidge\% against the ridge (DM
  \rhldDmLassoRidge).
\item The best classical extension, the OLS on \rhabBestOlsBench, changes QLIKE by \rhabOlsSharqlAllPct\% (DM
  \rhabOlsSharqlAllDm), less than any full set (Table~\ref{tab:rh_ablation}, Panel C).
\item On MSE the ranking reverses: the full set has a higher every-bar MSE than its HAR-set fit in
  \rhcaNMseHigherBarSpan{} of \rhorVsOlsFullN{} classes, by \rhcaLinMsePctBarSpanRng\% (linear) and
  \rhcaTreeMsePctBarSpanRng\% (trees), and \rhcaNMseHigherSigBarSpan{} of these increases are significant.
\item In \rhcaYearsOos, the \rhcaTopDays{} days with the largest squared-error difference, all between
  \rhcaTopFirstBarOos{} and \rhcaTopLastBarOos, hold \rhcaLinTopShareBarOosRng\% (linear) and
  \rhcaTreeTopShareBarOosRng\% (trees) of the full set's excess squared error.
\item Without the crash window from \rhcaCrashStart{} to \rhcaCrashEnd, the full set changes the every-bar MSE by
  \rhcaLinMsePctBarOosExCrashRng\% (linear) and \rhcaTreeMsePctBarOosExCrashRng\% (trees), significantly in
  \rhcaNMseLowerSigBarOosExCrash{} of \rhorVsOlsFullN{} classes.
\item By subperiod, the linear full sets raise the out-of-sample $R^2$ in \rhorSubALabel{} (\rhorSubARoosAfMin--\rhorSubARoosAfMax\%
  against \rhorSubARoosHar\%) and lower it in \rhorSubBLabel{} (\rhorSubBRoosAfMin--\rhorSubBRoosAfMax\% against
  \rhorSubBRoosHar\%).
\item The gain has two channels (Figure~\ref{fig:rh_channels}): the FOMC statement days hold \rhchEventBarPct\% of the
  bars and \rhchLinShareEventMin--\rhchLinShareEventMax\% (linear) and \rhchTreeShareEventMin--\rhchTreeShareEventMax\%
  (trees) of the gain, most of it on their bars from the statement to the close.
\item On statement days the FOMC group added alone to the HAR set delivers \rhchEvFomcShareLinRng\% (linear) and
  \rhchEvFomcShareTreeRng\% (trees) of the full set's gain.
\item The lasso wins on statement days and loses on ordinary days: against the ridge its QLIKE changes by
  \rhchLassoRidgeEventPct\% on statement days and \rhchLassoRidgeOrdPct\% on ordinary days ($p = \rhchLassoRidgeOrdP$).
\item On ordinary days the full set still gains, by \rhchCorrLinRng\% (linear) and \rhchCorrTreeRng\% (trees) under
  each way of removing the FOMC channel, and the ways agree within \rhchAgreePp{} percentage points for
  \rhchNCorrAgree{} classes.
\item On ordinary days the realized moments added alone give \rhchOrdRmShareTreeRng\% of the trees' gain and
  \rhchXEnetOrdRmVsHarShare--\rhchXLassoOrdRmVsHarShare\% of the linear models' gain.
\item The leave-one-out ablations agree: removing the realized moments raises QLIKE by
  \rhabLoRmAllPctMin--\rhabLoRmAllPctMax\%, removing the FOMC group by \rhabLoFomcAllPctMin--\rhabLoFomcAllPctMax\%, and
  removing any other group changes it by at most \rhabLoOtherAllMaxAbs\% (Table~\ref{tab:rh_ablation}, Panel A; the
  trees' arms other than FOMC are on the unpatched design).
\item An unpenalized OLS on the HAR set, the realized moments and the FOMC calendar changes QLIKE by
  \rhabOlsAddRmFomcVsHarAllPct\% against the OLS-HAR (DM \rhabOlsAddRmFomcVsHarAllDm), about as much as the penalized
  linear full sets.
\item The gain holds across windows: across \rhwsNWinWord{} windows from \rhwsWinMin{} to \rhwsWinMax{} days,
  the linear full sets beat the HAR set significantly in \rhwsBarNSigAll{} of \rhwsNCells{} cells, with
  QLIKE changes from \rhwsBarPctMin\% to \rhwsBarPctMax\%, and the trees in \rhlcNGapSigBar{} of \rhlcNGapCellsBar{} cells
  (Appendix~\ref{app:robust}).
\end{enumerate}

\input{generated/rehaul_ablation_table}

\input{generated/rehaul_channels_fig}

% ======================================================================================================================
\section{Beyond HAR, realized moments and the FOMC calendar}
\label{sec:rh_beyond}

\begin{enumerate}
\item The base set, the HAR set plus the realized moments and the FOMC calendar, reaches
  \rhabBaseShareMin--\rhabBaseShareMax\% of each class's full-set gain over its HAR fit.
\item Beyond the base, the penalized linear full sets are not significantly more accurate than the same class on the
  base, with smallest $p$-values \rhabLinFullVsAddRmFomcAllPMin{} at every bar and \rhabLinFullVsAddRmFomcOrdPMin{} on
  ordinary days (Table~\ref{tab:rh_ablation_beyond}).
% CONFLICT: under D2 the OLS is the primary benchmark, but the beyond-base comparison sets each class against its own
% penalized base (the full set's frozen penalties). Against the OLS on the base the elastic-net full set is better on
% ordinary days and worse at every bar (next beat), so "the linear models gain nothing beyond the base" depends on the
% base's estimator.
\item Against the OLS on the base, the elastic-net full set changes QLIKE by \rhbyFLFullOrdPct\% on ordinary days
  ($p = \rhbyFLFullOrdP$) and by \rhbyFLFullAllPct\% at every bar ($p = \rhbyFLFullAllP$), because the OLS base is
  better on statement days.
\item The boosted trees gain beyond the base: QLIKE changes by \rhabLgbmFullVsAddRmFomcAllPct\% for LightGBM (DM
  \rhabLgbmFullVsAddRmFomcAllDm) and \rhabXgbFullVsAddRmFomcAllPct\% for XGBoost (DM \rhabXgbFullVsAddRmFomcAllDm), and
  by \rhabRfFullVsAddRmFomcAllPct\% for the forest ($p = \rhabRfFullVsAddRmFomcAllP$).
\item The trees' gain beyond the base is on ordinary days, \rhabLgbmFullVsAddRmFomcOrdPct\% and
  \rhabXgbFullVsAddRmFomcOrdPct\%, and not significant on statement days ($p = \rhabLgbmFullVsAddRmFomcEvP$ and
  \rhabXgbFullVsAddRmFomcEvP).
\item The remaining groups overlap: added one at a time to the base, their gains sum to \rhabLgbmBeyondRatioAll{}
  (LightGBM) and \rhabXgbBeyondRatioAll{} (XGBoost) times the full set's gain.
\item The gain is not a few shared factors: the best model with at most $k = 3$ principal or partial-least-squares
  components of the remaining series recovers \rhbyFLgbmBestShortOrdR{} (LightGBM) and \rhbyFXgbBestShortOrdR{}
  (XGBoost) of the ordinary-day gain, and the pre-set reading for both is ``\rhbyFLgbmReading''.
\item The first $k = 3$ principal components hold \rhbyFLShareCumThree\% of the remaining series' variance, and the
  linear version of the test gives the same reading (``\rhbyFLReading'').
\item Ranked by importance, the single top remaining series recovers \rhbyLLgbmKOneAllR{} of LightGBM's gain beyond
  the base, but the full set stays more accurate (DM \rhbyLLgbmKOneAllDmFull), and the two boosted trees read
  differently (LightGBM \rhbyLLgbmReading, XGBoost \rhbyLXgbReading).
% CONFLICT: the JFEc budget puts the feature-count ladder figure in this section, but the generated figure
% (fig:rh_ladder) has a 16:00-bar panel, which D16 keeps out of the main text, and its generator places it in the
% appendix. The figure is not \input here; the next beat gives the every-bar rungs from macros.
\item On the full feature-count ladder, the trees refitted on their top $k = 10$ of \rhldNSeries{} series keep
  \rhldLgbmBarRTen\%, \rhldXgbBarRTen\% and \rhldRfBarRTen\% of the full set's every-bar gain, against
  \rhldLgbmBarRThree\%, \rhldXgbBarRThree\% and \rhldRfBarRThree\% at $k = 3$; the linear models at $k = 5$ keep
  \rhldRidgeBarRFive\%, \rhldLassoBarRFive\% and \rhldEnetBarRFive\%.
% CONFLICT: the SHAP generator splits the predictors into standard (HAR, calendar, the realized moments, the VIX, the
% futures volume and trade count) and non-standard series (the FOMC series among them), which does not match D18's base
% (HAR + RM + FOMC). No macro gives the attribution share of the six remaining groups, hence the pending mark below.
\item In the trees' attribution the HAR components hold \rhfipHarMin--\rhfipHarMax\% and the realized moments
  \rhfipMomMin--\rhfipMomMax\% of the mean absolute SHAP value \citep{lundberg2018consistent}, and the six remaining
  groups \rhPending{TreeSHAP share of the six groups beyond HAR, RM and FOMC} (Table~\ref{tab:rh_fip_shares}).
\item Use is not accuracy: the trees split on \rhfipLgbmNUsedAll, \rhfipXgbNUsedAll{} and \rhfipRfNUsedAll{} of the
  \rhfipNExog{} exogenous series at each of the \rhfipNOrigins{} origins, yet the remaining series change their QLIKE by
  \rhabLgbmFullVsAddRmFomcAllPct\% (LightGBM) and \rhabXgbFullVsAddRmFomcAllPct\% (XGBoost).
\item The linear models put more of their exogenous weight on the non-standard series than the trees do, in
  \rhfipNDLess{} of \rhfipNDAll{} tree-minus-linear comparisons.
\item The late-starting series gain attribution once they are live: the StockTwits share rises from
  \rhfilStwPreMin--\rhfilStwPreMax\% in \rhfilStwPreYears{} to \rhfilStwLiveMin--\rhfilStwLiveMax\% in
  \rhfilStwLiveYears, and the volatility-demand share from \rhfilVdPreMin--\rhfilVdPreMax\% to
  \rhfilVdLiveMin--\rhfilVdLiveMax\% in \rhfilVdLiveYears.
\end{enumerate}

\input{generated/rehaul_ablation_beyond_table}

\input{generated/rehaul_fip_shares}

% ======================================================================================================================
\section{Economic value: \zdte{} straddles}
\label{sec:rh_econ}\label{sec:rh_econ_measures}

% CONFLICT: D16 keeps the application on the forecast of the 16:00 bar ("keep the econ application as is") and keeps
% 16:00-bar accuracy out of the main text. A main-text reader therefore cannot see how accurate the forecast that signs
% the trade is (the linear full sets are not more accurate than the OLS-HAR at that bar; Appendix). One sentence
% pointing to the appendix is proposed in the first beat.
\begin{enumerate}
\item On each option day the strategy trades the at-the-money SPXW straddle that expires that day, half an hour before
  the close, long when the pooled forecast of the closing bar's realized variance exceeds the straddle's implied
  variance and short otherwise; the accuracy of that bar's forecast is in Appendix~\ref{app:rh_closing_bench}.
\item The benchmarks are the same strategy signed by the OLS-HAR forecast and the short straddle held every day.
\item The statistics use the \rheSzDays{} option days after a warm-up of \rheSzWarmup{} days, of the \rheNdays{} option
  days from \rheSampleStart{} to \rheSampleEnd.
\item The Sharpe ratio pools the daily returns, each day weighted equally, and annualizes by the square root of
  \rheAnnDays{} option days per calendar year; differences are tested with \citet{LedoitWolf2008} (block length
  \rheLWBlock, \rheLWBoot{} resamples).
\item The full-set strategies have a median Sharpe ratio of \rhemFullShMed{} at mid-quotes and \rhemFullShXMed{} at
  quoted bid and ask, against \rhemShOlsHar{} and \rhemShXOlsHar{} for the OLS-HAR and \rhemShShort{} and \rhemShXShort{}
  for the short straddle (Table~\ref{tab:rh_econ_main}).
\item Against the OLS-HAR strategy, the full-set Sharpe ratios are higher in \rhemFullVsOlsPos{} of \rhmcsNAlt{}
  classes, by \rhemFullVsOlsMin{} to \rhemFullVsOlsMax, significantly for \rhemFullVsOlsSigWord{} (the random forest,
  $p = \rhemPDShRfFullVsOls$).
\item Against the short straddle, \rhemFullVsShortSigWord{} full-set strategies have a significantly higher Sharpe
  ratio, against \rhemHarVsShortSigWord{} of the HAR-set strategies.
\item Within each class the full set raises the Sharpe ratio over the class's HAR set in \rhemFullVsHarPosWord{}
  classes, significantly in \rhemFullVsHarSigWord.
\item By calendar year the median full-set Sharpe ratio is above the median HAR-set ratio in
  \rhemYearsFullMedAboveHarMedWord{} of \rhemNYearsWord{} years, and the \rhmcsNAltWord{} full-set ratios are positive in
  \rhemYearsFullAllPosWord{} of \rhemNYearsWord{} years (Figure~\ref{fig:rh_econ_byyear}).
\item The profits come from the long days: the full-set strategies go long on \rhedLongPctMin--\rhedLongPctMax\% of
  the days and earn \rhedLongShareMin--\rhedLongShareMax\% of their summed return on them, while the HAR-set strategies
  go long on \rhedHarLongNMin--\rhedHarLongNMax{} of the \rhedNDays{} days.
\item The profits come from a few days: the \rhedNTop{} best days hold \rhedBestShareMin--\rhedBestShareMax\% of the
  summed gain over the short straddle, and without them the Sharpe ratio difference against the short straddle is
  negative for \rhedDShortExBestNeg{} of \rhmcsNAlt{} strategies (\rhedDShortExBestMin{} to \rhedDShortExBestMax).
\item Without those days the full set still beats its HAR-set twin, by \rhedDHarExBestMin{} to \rhedDHarExBestMax{}
  in Sharpe ratio.
\item The full-set strategies break even at a cost of \rhemBeMeanFullMinPct--\rhemBeMeanFullMaxPct\% of the premium,
  against a median quoted half-spread of \rhemHalfSpreadMed\%.
\item The strategy returns carry no significant market beta (\rhebBetaSigFullWord{} of the full-set strategies), and
  their alphas are significant for \rhebAlphaSigFullWord{} (Appendix~\ref{app:econ}).
\end{enumerate}

\input{generated/rehaul_econ_main_table}

\input{generated/rehaul_econ_byyear_fig}

% ======================================================================================================================
\section{Conclusion}
\label{sec:concl}

\begin{enumerate}
\item Pooled over the regular-hours bars, \rhdmNSeriesPooled{} exogenous series lower the QLIKE of realized-variance
  forecasts against the OLS-HAR by \rhorVsOlsRedLinRng\% (linear) and \rhorVsOlsRedTreeRng\% (trees).
\item The same forecasts have a higher MSE, which a few crash days produce; without the crash window the full set
  also lowers the MSE.
\item The gain runs through the FOMC calendar on statement days and the futures' realized moments on ordinary days,
  which together reach \rhabBaseShareMin--\rhabBaseShareMax\% of it.
\item What the remaining \rhbyFLNSeries{} series add is not significant for penalized linear models and is a small gain
  for boosted trees, spread over the series rather than held in a few components.
\item A practitioner who wants most of the gain needs the HAR set, the futures' realized moments and the FOMC calendar;
  the other series pay only inside a tree ensemble.
\item The forecasts carry economic value in a \zdte{} straddle strategy, but the evidence rests on few days and is
  significant against the OLS-HAR strategy for \rhemFullVsOlsSigWord{} class.
\item The StockTwits and volatility-demand series are live on \rhlvStSentEvRange\% and \rhlvVdEvRange\% of the scored
  bars, so their contribution is measured on short windows.
\end{enumerate}

% ======================================================================================================================
\appendix
\section{Online appendix: pointers}
\label{app:pointers}

\subsection{Tuning disclosure}
\label{sec:rh_methods_hyper}
\begin{enumerate}
\item Re-tuning the penalties every year moves the linear full sets' every-bar QLIKE change to
  between \rhtuAnnPctBarMin\% and \rhtuAnnPctBarMax\%, and the neighbouring penalties give changes between
  \rhtuNbrPctBarMin\% and \rhtuNbrPctBarMax\%.
\item The forest runs on the oneDAL engine, \rhtuRfSpeedup{} times faster than the stock forest, with a validation
  QLIKE within \rhtuRfQlikeDiffPct\% of it.
\end{enumerate}

\subsection{The closing half hour}
\label{app:rh_closing_bench}\label{app:rh_closing}
% CONFLICT: D3 corrals every closing-interval model (CI-HAR, close-trained linear and trees) into a final main-text
% section, while D16 and the JFEc budget keep 16:00-bar accuracy in the appendix. Here they are appendix only.
\begin{enumerate}
\item The accuracy at the closing bar, of the pooled models and of the models trained on the last half hour only
  (Table~\ref{tab:rh_oos_last30}), the rebuilt closing-interval HAR and the pre-registered three-way training
  comparison go here.
\end{enumerate}

\input{generated/rehaul_results_oos_last30}

\subsection{Robustness}
\label{app:robust}
\begin{enumerate}
\item The windows, the training arms, the holiday sessions, the model confidence set, SPA and StepM tests go here.
\end{enumerate}

\subsection{Economics details}
\label{app:econ}
\begin{enumerate}
\item The Sharpe ratio definitions, sizing, compounding, the interval-shift test, break-evens, the market beta and the
  allocation with the market, and the profit decomposition go here.
\end{enumerate}

\bibliographystyle{plainnat}
\bibliography{../references}

\end{document}
